\documentclass[11pt]{article}
\usepackage[a4paper, margin=1in]{geometry}
\usepackage{hyperref}
\usepackage{titlesec}
\usepackage{enumitem}
\usepackage{listings}
\usepackage{xcolor}
\usepackage{tcolorbox}
\usepackage{amsmath}
\usepackage{graphicx}

\definecolor{codegray}{gray}{0.95}
\definecolor{darkblue}{rgb}{0.0,0.0,0.5}
\definecolor{explainblue}{rgb}{0.0,0.3,0.6}

\lstset{
  backgroundcolor=\color{codegray},
  basicstyle=\ttfamily\small,
  breaklines=true,
  frame=single,
  keywordstyle=\color{darkblue}\bfseries,
  commentstyle=\itshape\color{gray},
  stringstyle=\color{black}
}

\newtcolorbox{noteBox}{
  colback=blue!5,
  colframe=blue!60,
  title=\textbf{💡 Note Importante},
  sharp corners,
  fonttitle=\bfseries
}

\newtcolorbox{warnBox}{
  colback=red!5,
  colframe=red!60,
  title=\textbf{⚠️ Attention},
  sharp corners,
  fonttitle=\bfseries
}

\newtcolorbox{exampleBox}{
  colback=green!5,
  colframe=green!60,
  title=\textbf{📝 Exemple Concret},
  sharp corners,
  fonttitle=\bfseries
}

\newtcolorbox{explainBox}{
  colback=yellow!5,
  colframe=orange!60,
  title=\textbf{🔍 Explication Détaillée},
  sharp corners,
  fonttitle=\bfseries
}

\title{\textbf{Cours Complet de Révision - Sui Move}\\
\large De Zéro à Confiant : Guide Explicatif Détaillé\\
\large Abilities, Objects, OTW, Publisher, Capabilities, Option, Vectors, Loops, Events}
\author{Bootcamp Notes - Version Explicative}
\date{}

\begin{document}

\maketitle
\tableofcontents
\newpage

% ============================================================
\section{Introduction : Comment Utiliser Ce Document}

\begin{noteBox}
Ce document est conçu pour la \textbf{révision complète}. Chaque concept est expliqué en détail avec des analogies simples, pas seulement du code. Lisez les explications avant de regarder le code.
\end{noteBox}

\subsection{Structure du Document}
\begin{itemize}
  \item \textbf{Explications détaillées} : Pourquoi chaque concept existe
  \item \textbf{Analogies simples} : Comparaisons avec la vie réelle
  \item \textbf{Exemples concrets} : Code commenté ligne par ligne
  \item \textbf{Points clés} : Ce qu'il faut retenir absolument
\end{itemize}

% ============================================================
\section{Les Bases : Move et Sui (Comprendre le Fondement)}

\subsection{Qu'est-ce qu'un Smart Contract ? (Explication Simple)}

Imaginez un \textbf{contrat intelligent} comme un \textbf{automate de confiance}. 

\begin{explainBox}
\textbf{Analogie} : C'est comme une machine à café qui :
\begin{itemize}
  \item Accepte toujours les mêmes pièces (déterministe)
  \item Donne toujours le même café pour les mêmes pièces (vérifiable)
  \item Ne peut pas être triché - elle suit strictement ses règles (sécurisé)
\end{itemize}
\end{explainBox}

\subsubsection{Les 3 Propriétés Essentielles}

\begin{enumerate}
  \item \textbf{Déterministe} : Même entrée = même résultat, toujours
  \begin{itemize}
    \item Si vous donnez 2€ à la machine, vous obtenez toujours le même café
    \item Le résultat ne dépend pas de l'heure, du jour, ou de qui utilise la machine
  \end{itemize}
  
  \item \textbf{Vérifiable} : N'importe qui peut vérifier que le résultat est correct
  \begin{itemize}
    \item Les validateurs du réseau peuvent rejouer votre transaction et obtenir le même résultat
    \item C'est comme avoir des témoins qui vérifient que la machine a bien fonctionné
  \end{itemize}
  
  \item \textbf{Sécurisé} : Impossible de dupliquer des actifs ou de tricher
  \begin{itemize}
    \item La machine ne peut pas créer de l'argent à partir de rien
    \item Une fois qu'un café est donné, il ne peut pas être "dupliqué"
  \end{itemize}
\end{enumerate}

\subsection{Qu'est-ce que Move ? (Le Langage pour les Actifs)}

Move n'est pas juste un langage de programmation - c'est un langage \textbf{spécialement conçu pour gérer des actifs numériques}.

\begin{explainBox}
\textbf{La Philosophie de Move} :
\begin{itemize}
  \item Dans la plupart des langages, un nombre est juste un nombre : \texttt{balance = 100}
  \item Dans Move, un actif est \textbf{spécial} : il ne peut pas être copié, perdu, ou créé par accident
  \item Move force le développeur à être explicite sur ce qu'il fait avec les actifs
\end{itemize}
\end{explainBox}

\subsubsection{Pourquoi Move est Différent}

\begin{exampleBox}
\textbf{Exemple de Problème Résolu} :

Dans un langage classique, vous pourriez accidentellement faire :
\begin{lstlisting}
let my_coin = 100;
let your_coin = my_coin;  // Maintenant vous avez tous les deux 100 !
\end{lstlisting}

Dans Move, cela est \textbf{impossible} pour les actifs. Si vous transférez un actif, il est \textbf{déplacé}, pas copié.
\end{exampleBox}

\subsection{Qu'est-ce que Sui Move ? (Move + Système d'Objets)}

Sui Move = Move + Système d'objets unique de Sui

\begin{explainBox}
\textbf{La Différence Clé} :
\begin{itemize}
  \item \textbf{Move classique} : Gère les actifs de manière sécurisée
  \item \textbf{Sui Move} : Ajoute le concept d'\textbf{objets} avec propriétaires, IDs uniques, et transferts rapides
\end{itemize}
\end{explainBox}

\subsubsection{Le Concept d'Objet dans Sui}

Dans Sui, presque tout ce qui compte est un \textbf{objet} :
\begin{itemize}
  \item Un NFT est un objet
  \item Une pièce de monnaie est un objet
  \item Un ticket d'événement est un objet
  \item Même les "capabilities" (permissions) sont des objets
\end{itemize}

\begin{noteBox}
\textbf{Pensez-y ainsi} : Dans Sui, au lieu de dire "j'ai 100 tokens dans mon compte", vous dites "je possède 100 objets de type Token". Chaque objet a un ID unique et peut être transféré individuellement.
\end{noteBox}

\subsection{Modules, Packages et Adresses (L'Organisation du Code)}

\begin{explainBox}
\textbf{Analogie avec un Immeuble} :
\begin{itemize}
  \item \textbf{Package} = L'immeuble entier (ce que vous publiez sur Sui)
  \item \textbf{Module} = Un appartement dans l'immeuble (un fichier de code)
  \item \textbf{Address} = L'adresse de l'immeuble sur la blockchain
\end{itemize}
\end{explainBox}

\subsubsection{Explication Détaillée}

\begin{enumerate}
  \item \textbf{Package} : Collection de modules liés
  \begin{itemize}
    \item C'est ce que vous publiez sur Sui en une seule fois
    \item Exemple : Un package "NFT\_Collection" pourrait contenir plusieurs modules
  \end{itemize}
  
  \item \textbf{Module} : Un fichier de code avec des fonctions
  \begin{itemize}
    \item Format : \texttt{package\_name::module\_name}
    \item Exemple : \texttt{my\_nft::minting} (module "minting" dans le package "my\_nft")
  \end{itemize}
  
  \item \textbf{Address} : L'identifiant unique du package sur la blockchain
  \begin{itemize}
    \item Chaque package a une adresse unique (son Package ID)
    \item C'est comme l'adresse GPS de votre immeuble
  \end{itemize}
\end{enumerate}

% ============================================================
\section{Les Types et Structs (Les Briques de Base)}

\subsection{Les Types de Base (Les Outils Fondamentaux)}

\begin{explainBox}
Les types de base sont comme les \textbf{outils de base} dans une boîte à outils. Vous les combinez pour créer des choses plus complexes.
\end{explainBox}

\subsubsection{Les Types Principaux}

\begin{itemize}
  \item \texttt{u64} : Un nombre entier positif (0, 1, 2, ... jusqu'à très grand)
  \begin{itemize}
    \item Le "u" signifie "unsigned" (non signé) - donc pas de nombres négatifs
    \item Le "64" signifie 64 bits - peut stocker des nombres très grands
  \end{itemize}
  
  \item \texttt{bool} : Vrai ou Faux (true/false)
  \begin{itemize}
    \item Comme une question à réponse oui/non
    \item Exemple : "Est-ce que l'utilisateur est admin ?" → true ou false
  \end{itemize}
  
  \item \texttt{address} : Une adresse sur la blockchain
  \begin{itemize}
    \item Comme une adresse postale, mais pour la blockchain
    \item Identifie de manière unique un compte ou un contrat
  \end{itemize}
\end{itemize}

\subsection{Les Structs (Créer Vos Propres Types)}

Un \textbf{struct} est comme créer un \textbf{modèle personnalisé} pour vos données.

\begin{explainBox}
\textbf{Analogie avec un Formulaire} :
\begin{itemize}
  \item Un struct est comme un formulaire avec des champs à remplir
  \item Chaque champ a un nom et un type
  \item Une fois rempli, vous avez une "instance" du struct
\end{itemize}
\end{explainBox}

\subsubsection{Exemple Simple : Un Point 2D}

\begin{lstlisting}[language=Rust, caption=Struct Point - Exemple Basique]
module demo::basic {
    // On définit un nouveau type appelé "Point"
    struct Point {
        x: u64,  // Coordonnée X (un nombre)
        y: u64   // Coordonnée Y (un nombre)
    }
    
    // Maintenant on peut créer des Points :
    // Point { x: 10, y: 20 } crée un point à la position (10, 20)
}
\end{lstlisting}

\begin{noteBox}
\textbf{Points Clés} :
\begin{itemize}
  \item Un struct est juste une \textbf{définition} - comme un plan
  \item Pour l'utiliser, vous devez créer une \textbf{instance} avec des valeurs
  \item Les structs sont la façon principale de représenter des données dans Move
\end{itemize}
\end{noteBox}

% ============================================================
\section{Les Abilities (Extrêmement Important - Comprendre les Permissions)}

\subsection{Qu'est-ce qu'une Ability ? (Le Concept Fondamental)}

Une \textbf{ability} est une \textbf{permission} sur un type. Elle contrôle ce que vous pouvez faire avec les valeurs de ce type.

\begin{explainBox}
\textbf{Analogie avec les Permissions de Fichier} :
\begin{itemize}
  \item Sur votre ordinateur, un fichier peut être "lecture seule" ou "modifiable"
  \item Dans Move, un type peut avoir différentes "abilities" qui contrôlent ce qu'on peut en faire
  \item C'est comme donner des permissions spécifiques à chaque type de données
\end{itemize}
\end{explainBox}

\subsection{Les 4 Abilities (Les Permissions Possibles)}

\subsubsection{1. \texttt{copy} - La Permission de Copier}

\begin{explainBox}
\textbf{Qu'est-ce que copy ?}
\begin{itemize}
  \item Permet de \textbf{dupliquer} une valeur
  \item Après copie, vous avez deux copies identiques
  \item \textbf{DANGER} : Si un actif a \texttt{copy}, vous pouvez le dupliquer = créer de l'argent à partir de rien !
\end{itemize}
\end{explainBox}

\begin{warnBox}
\textbf{⚠️ Règle d'Or} : Les actifs (argent, NFTs, tokens) ne doivent \textbf{JAMAIS} avoir \texttt{copy}. Sinon, vous pouvez dupliquer de l'argent, ce qui détruit toute la sécurité.
\end{warnBox}

\subsubsection{2. \texttt{drop} - La Permission de Jeter}

\begin{explainBox}
\textbf{Qu'est-ce que drop ?}
\begin{itemize}
  \item Permet de \textbf{détruire} ou \textbf{ignorer} une valeur
  \item Si une valeur n'a pas \texttt{drop}, vous devez l'utiliser quelque part (la déplacer, la transférer, etc.)
  \item C'est comme pouvoir jeter quelque chose à la poubelle
\end{itemize}
\end{explainBox}

\begin{exampleBox}
\textbf{Exemple} :
\begin{itemize}
  \item Un \texttt{Point} peut avoir \texttt{drop} - si vous ne l'utilisez pas, vous pouvez l'ignorer
  \item Un \texttt{Coin} (pièce) ne devrait \textbf{pas} avoir \texttt{drop} - vous ne pouvez pas juste "jeter" de l'argent !
\end{itemize}
\end{exampleBox}

\subsubsection{3. \texttt{store} - La Permission de Stocker}

\begin{explainBox}
\textbf{Qu'est-ce que store ?}
\begin{itemize}
  \item Permet de \textbf{stocker} la valeur à l'intérieur d'autres structs
  \item Si un struct n'a pas \texttt{store}, vous ne pouvez pas le mettre dans un autre struct
  \item C'est comme pouvoir mettre quelque chose dans une boîte
\end{itemize}
\end{explainBox}

\begin{noteBox}
\textbf{Exemple Concret} :
\begin{itemize}
  \item Si vous voulez créer un struct \texttt{Wallet} qui contient des \texttt{Coin}, alors \texttt{Coin} doit avoir \texttt{store}
  \item Sinon, vous ne pouvez pas mettre le \texttt{Coin} dans le \texttt{Wallet}
\end{itemize}
\end{noteBox}

\subsubsection{4. \texttt{key} - La Permission d'Être un Objet Sui}

\begin{explainBox}
\textbf{Qu'est-ce que key ?}
\begin{itemize}
  \item Permet à un struct d'exister comme un \textbf{objet top-level} sur Sui
  \item Un objet avec \texttt{key} a un ID unique et peut être possédé ou partagé
  \item C'est la permission spéciale qui fait qu'un struct devient un "objet Sui"
\end{itemize}
\end{explainBox}

\begin{noteBox}
\textbf{Règle Importante} : Si un struct a \texttt{key}, il \textbf{DOIT} contenir un champ \texttt{id: UID}. C'est l'identité unique de l'objet.
\end{noteBox}

\subsection{Pourquoi les Abilities Sont Cruciales (Exemples Concrets)}

\subsubsection{Exemple 1 : Données Sûres Peuvent Avoir Copy}

\begin{lstlisting}[language=Rust, caption=Point avec Copy - C'est OK]
module demo::abilities {
    // Un Point est juste des coordonnées - pas un actif
    // Donc on peut le copier sans danger
    struct Point has copy, drop, store {
        x: u64,
        y: u64
    }
    
    // Maintenant vous pouvez faire :
    // let p1 = Point { x: 10, y: 20 };
    // let p2 = p1;  // ✅ OK - copie autorisée
    // Vous avez maintenant deux Points identiques
}
\end{lstlisting}

\begin{noteBox}
\textbf{Pourquoi c'est sûr ?} : Copier des coordonnées ne crée pas de valeur - c'est juste des données. Personne ne peut "gagner" quelque chose en copiant un Point.
\end{noteBox}

\subsubsection{Exemple 2 : Les Actifs NE Doivent PAS Avoir Copy}

\begin{lstlisting}[language=Rust, caption=Coin SANS Copy - C'est Correct]
module demo::abilities {
    // Un Coin représente de l'argent
    // On NE veut PAS copy - sinon on peut dupliquer l'argent !
    struct Coin has store {
        value: u64
    }
    
    // Maintenant vous NE pouvez PAS faire :
    // let coin1 = Coin { value: 100 };
    // let coin2 = coin1;  // ❌ ERREUR - copy interdit !
    
    // Vous DEVEZ déplacer (move) le coin :
    // transfer::transfer(coin1, recipient);  // ✅ OK - déplacement
}
\end{lstlisting}

\begin{warnBox}
\textbf{⚠️ Erreur Critique} : Si vous ajoutez \texttt{copy} à un struct qui représente un actif (comme \texttt{Coin}), vous créez une faille de sécurité majeure. N'importe qui pourrait dupliquer de l'argent !
\end{warnBox}

\subsection{La Règle des "Field Abilities" (Règle de Composition)}

\begin{explainBox}
\textbf{La Règle} : Un struct ne peut avoir une ability que si \textbf{tous ses champs} supportent cette ability.
\end{explainBox}

\subsubsection{Exemple Concret}

\begin{lstlisting}[language=Rust, caption=Exemple de Field Ability Rule]
module demo::abilities {
    struct Coin has store {
        value: u64  // u64 a copy, drop, store
    }
    // ✅ Coin peut avoir store car u64 a store
    
    struct Wallet {
        coin: Coin  // Coin a store
    }
    // ❌ Wallet ne peut PAS avoir store car Coin n'a pas copy
    // (Pour avoir store, tous les champs doivent avoir copy OU store)
}
\end{lstlisting}

\begin{noteBox}
\textbf{Pensez-y ainsi} : Si vous voulez mettre quelque chose dans une boîte (\texttt{store}), tous les éléments à l'intérieur doivent pouvoir être mis dans une boîte aussi.
\end{noteBox}

% ============================================================
\section{Les Objets Sui et UID (Le Cœur du Système Sui)}

\subsection{Qu'est-ce qu'un Objet dans Sui ? (Concept Fondamental)}

Un \textbf{objet Sui} est une pièce de données qui a :
\begin{enumerate}
  \item Un \textbf{ID unique} (comme un numéro de série)
  \item Un \textbf{propriétaire} (ou est partagé)
  \item Des \textbf{champs de données}
\end{enumerate}

\begin{explainBox}
\textbf{Analogie avec un Objet Physique} :
\begin{itemize}
  \item Imaginez un \textbf{livre physique} :
  \begin{itemize}
    \item Il a un numéro ISBN unique (l'ID)
    \item Il appartient à quelqu'un (le propriétaire)
    \item Il contient des pages avec du texte (les données)
  \end{itemize}
  \item Un objet Sui fonctionne de la même manière, mais numériquement
\end{itemize}
\end{explainBox}

\subsection{Qu'est-ce que UID ? (L'Identité Unique)}

\texttt{UID} signifie \textbf{Unique Identifier} (Identifiant Unique).

\begin{explainBox}
\textbf{Qu'est-ce que UID ?}
\begin{itemize}
  \item C'est un type spécial fourni par Sui
  \item Il garantit que chaque objet a un ID \textbf{vraiment unique}
  \item C'est comme un numéro de passeport - chaque personne en a un unique
\end{itemize}
\end{explainBox}

\subsubsection{Structure d'un Objet Sui}

\begin{lstlisting}[language=Rust, caption=Structure Basique d'un Objet Sui]
module demo::objects {
    use sui::object::{UID};
    
    // Pour qu'un struct soit un objet Sui, il DOIT :
    // 1. Avoir l'ability "key"
    // 2. Contenir un champ "id: UID"
    struct MyObject has key {
        id: UID,      // ✅ Obligatoire - l'identité unique
        value: u64    // Vos données personnalisées
    }
}
\end{lstlisting}

\begin{noteBox}
\textbf{Règle Absolue} : Si un struct a \texttt{key}, il \textbf{DOIT} contenir \texttt{id: UID}. C'est la seule façon pour Sui de suivre et gérer l'objet.
\end{noteBox}

\subsection{Objets Possédés vs Objets Partagés (Performance)}

\begin{explainBox}
\textbf{La Différence Cruciale} :
\begin{itemize}
  \item \textbf{Objet possédé} : Appartient à une seule adresse
  \begin{itemize}
    \item Transactions généralement \textbf{rapides} (pas besoin d'ordre global)
    \item Comme un livre dans votre bibliothèque personnelle
  \end{itemize}
  
  \item \textbf{Objet partagé} : Accessible par tout le monde
  \begin{itemize}
    \item Transactions \textbf{plus lentes} (nécessite un ordre de consensus)
    \item Comme un livre dans une bibliothèque publique - tout le monde peut l'emprunter
  \end{itemize}
\end{itemize}
\end{explainBox}

\subsubsection{Quand Utiliser Chacun ?}

\begin{exampleBox}
\textbf{Objets Possédés} - Utilisez pour :
\begin{itemize}
  \item NFTs personnels
  \item Pièces de monnaie dans un portefeuille
  \item Tickets d'événement
  \item Tout ce qui appartient à une seule personne
\end{itemize}

\textbf{Objets Partagés} - Utilisez pour :
\begin{itemize}
  \item Configuration globale d'une application
  \item Registres accessibles par tous
  \item Systèmes de gouvernance
  \item Tout ce qui doit être accessible par plusieurs personnes
\end{itemize}
\end{exampleBox}

\begin{warnBox}
\textbf{⚠️ Attention Performance} : Si vous mettez un énorme vecteur dans un objet partagé et que beaucoup d'utilisateurs le modifient, les performances peuvent se dégrader (beaucoup de conflits).
\end{warnBox}

% ============================================================
\section{Publisher (Preuve d'Auteur de Package)}

\subsection{Qu'est-ce que Publisher ? (Le Concept)}

\textbf{Publisher} est un objet spécial créé automatiquement quand vous publiez un package sur Sui.

\begin{explainBox}
\textbf{Analogie avec un Certificat d'Auteur} :
\begin{itemize}
  \item Quand vous publiez un livre, vous recevez un certificat qui dit "Je suis l'auteur de ce livre"
  \item \textbf{Publisher} est ce certificat, mais pour votre code
  \item Il prouve : "Je suis l'auteur original de ce package"
\end{itemize}
\end{explainBox}

\subsection{Pourquoi Publisher Existe (Le Problème Résolu)}

\begin{explainBox}
\textbf{Le Problème} : Comment s'assurer que seuls les créateurs originaux peuvent faire des actions administratives ?

\textbf{La Solution} : Publisher est la preuve on-chain que vous êtes le créateur original.
\end{explainBox}

\subsubsection{Exemples d'Utilisation}

\begin{exampleBox}
\textbf{Cas d'Usage Typiques} :
\begin{itemize}
  \item \textbf{Mintage restreint} : Seul le créateur peut mint de nouveaux NFTs
  \item \textbf{Mise à jour de règles} : Seul le créateur peut changer les règles du contrat
  \item \textbf{Upgrade de package} : Selon les paramètres, seul le créateur peut mettre à jour le code
\end{itemize}
\end{exampleBox}

\subsection{Comment Publisher est Utilisé (Pattern Commun)}

\begin{explainBox}
\textbf{Le Pattern} :
\begin{enumerate}
  \item Votre fonction admin nécessite \texttt{\&Publisher} comme paramètre
  \item Si l'appelant ne possède pas le Publisher, il ne peut pas appeler la fonction
  \item Seul le créateur original (qui possède le Publisher) peut faire l'action
\end{enumerate}
\end{explainBox}

\begin{lstlisting}[language=Rust, caption=Exemple d'Utilisation de Publisher]
module demo::publisher_example {
    use sui::package::Publisher;
    
    // Seulement le créateur (qui a Publisher) peut mint
    public fun mint_nft(
        publisher: &Publisher,  // ✅ Doit passer Publisher
        ctx: &mut TxContext
    ) {
        // Logique de mint...
        // Si quelqu'un essaie d'appeler sans Publisher, ça échoue
    }
}
\end{lstlisting}

\begin{noteBox}
\textbf{Important} : Dans le code réel, vous ne vérifiez pas manuellement si Publisher est passé - le système Move le fait automatiquement. Si la fonction nécessite \texttt{\&Publisher} et que l'appelant ne l'a pas, la transaction échoue.
\end{noteBox}

% ============================================================
\section{One-Time Witness (OTW) - Initialiser Exactement Une Fois}

\subsection{Le Problème que OTW Résout}

Parfois, vous devez vous assurer qu'une initialisation se produit \textbf{exactement une fois}, jamais plus.

\begin{explainBox}
\textbf{Exemples de Problèmes Résolus} :
\begin{itemize}
  \item \textbf{Créer un type de token} : Vous ne voulez créer qu'un seul type "MyToken"
  \item \textbf{Configuration globale} : Vous ne voulez qu'une seule configuration partagée
  \item \textbf{Registre unique} : Vous ne voulez qu'un seul registre pour votre application
\end{itemize}
\end{explainBox}

\subsection{Qu'est-ce qu'un OTW ? (Le Mécanisme)}

Un \textbf{One-Time Witness} est un struct spécial qui :
\begin{enumerate}
  \item Est créé \textbf{une seule fois} (au moment de la publication)
  \item Est ensuite \textbf{consommé} (utilisé et détruit)
  \item Ne peut \textbf{jamais être recréé}
\end{enumerate}

\begin{explainBox}
\textbf{Analogie avec un Ticket Unique} :
\begin{itemize}
  \item Imaginez un ticket de concert spécial qui ne peut être utilisé qu'une fois
  \item Une fois utilisé, il est détruit et ne peut jamais être recréé
  \item OTW fonctionne de la même manière - c'est un "ticket" pour initialiser une fois
\end{itemize}
\end{explainBox}

\subsection{Pattern OTW Simple}

\begin{lstlisting}[language=Rust, caption=Exemple de Pattern OTW]
module demo::otw {
    // Le struct OTW - doit être vide et sans abilities
    struct INIT {}  // ✅ Vide, pas d'abilities = OTW parfait
    
    // Fonction d'initialisation - ne peut être appelée qu'une fois
    public fun init(
        witness: INIT,  // ✅ Le witness est consommé (détruit) après usage
        ctx: &mut TxContext
    ) {
        // Cette logique ne peut s'exécuter qu'une fois
        // car INIT ne peut être créé qu'une seule fois (à la publication)
        
        // Créer votre configuration globale, registre, etc.
    }
}
\end{lstlisting}

\begin{noteBox}
\textbf{Pourquoi ça marche ?} :
\begin{itemize}
  \item Le struct \texttt{INIT} est créé automatiquement une seule fois lors de la publication
  \item Une fois utilisé dans \texttt{init()}, il est détruit
  \item Il ne peut jamais être recréé - donc \texttt{init()} ne peut être appelée qu'une fois
\end{itemize}
\end{noteBox}

\begin{warnBox}
\textbf{⚠️ Erreur Commune} : Si vous donnez des abilities (comme \texttt{store} ou \texttt{copy}) au struct OTW, vous cassez la garantie "one-time". Le struct pourrait être stocké ou copié, permettant plusieurs initialisations.
\end{warnBox}

% ============================================================
\section{Capabilities (Permission comme Objet)}

\subsection{Qu'est-ce qu'une Capability ? (Le Concept)}

Une \textbf{capability} est un \textbf{objet} qui représente une permission.

\begin{explainBox}
\textbf{Changement de Paradigme} :
\begin{itemize}
  \item \textbf{Ancienne façon} : "Si l'expéditeur est admin, alors autoriser"
  \item \textbf{Façon Sui} : "Si vous possédez AdminCap, alors vous êtes admin"
\end{itemize}
\end{explainBox}

\subsubsection{Analogie avec une Clé Physique}

\begin{exampleBox}
\textbf{Analogie Simple} :
\begin{itemize}
  \item Une capability est comme une \textbf{clé physique}
  \item Si vous avez la clé de la porte, vous pouvez entrer
  \item Si vous n'avez pas la clé, vous ne pouvez pas entrer
  \item La clé peut être transférée à quelqu'un d'autre
\end{itemize}
\end{exampleBox}

\subsection{Pourquoi les Capabilities sont Puissantes}

\begin{explainBox}
\textbf{Avantages des Capabilities} :
\begin{enumerate}
  \item \textbf{Infalsifiables} : Seul votre module peut créer les capabilities
  \begin{itemize}
    \item C'est comme si vous étiez le seul à pouvoir fabriquer les clés
    \item Personne d'autre ne peut créer une fausse AdminCap
  \end{itemize}
  
  \item \textbf{Transférables} : Vous pouvez donner la capability à quelqu'un d'autre
  \begin{itemize}
    \item Comme donner votre clé à un ami
    \item Utile pour changer d'admin ou déléguer des permissions
  \end{itemize}
  
  \item \textbf{Stockables dans des Multisigs ou DAOs} : Les capabilities peuvent être gérées collectivement
  \begin{itemize}
    \item Plusieurs personnes peuvent posséder des capabilities
    \item Permet une gouvernance décentralisée
  \end{itemize}
\end{enumerate}
\end{explainBox}

\subsection{Exemple Débutant : MintCap}

\begin{lstlisting}[language=Rust, caption=Exemple Simple de Capability]
module demo::capabilities {
    use sui::object::{UID};
    use sui::tx_context::TxContext;
    
    // Un objet Coin (pièce de monnaie)
    struct Coin has key {
        id: UID,
        value: u64
    }
    
    // La Capability : Qui possède ceci a l'autorité de mint
    struct MintCap has key {
        id: UID
    }
    
    // Fonction de mint - nécessite MintCap
    public fun mint(
        cap: &MintCap,  // ✅ Doit passer MintCap
        amount: u64,
        ctx: &mut TxContext
    ): Coin {
        // Si l'utilisateur n'a pas MintCap, il ne peut pas appeler cette fonction
        Coin { 
            id: object::new(ctx), 
            value: amount 
        }
    }
}
\end{lstlisting}

\begin{noteBox}
\textbf{Comment ça fonctionne} :
\begin{itemize}
  \item L'utilisateur doit passer \texttt{\&MintCap} à la fonction \texttt{mint()}
  \item S'il ne possède pas \texttt{MintCap}, il ne peut pas le fournir
  \item Donc, seuls ceux qui possèdent \texttt{MintCap} peuvent mint
\end{itemize}
\end{noteBox}

\subsection{Capabilities vs Publisher vs OTW (Comparaison)}

\begin{explainBox}
\textbf{Les Différences Clés} :
\begin{itemize}
  \item \textbf{Publisher} : Preuve d'auteur de package (niveau package)
  \begin{itemize}
    \item Créé automatiquement à la publication
    \item Prouve que vous êtes l'auteur original
    \item Utilisé pour des actions au niveau du package
  \end{itemize}
  
  \item \textbf{OTW} : Garantit l'initialisation une seule fois
  \begin{itemize}
    \item Créé automatiquement une fois à la publication
    \item Consommé après usage
    \item Utilisé pour l'initialisation unique
  \end{itemize}
  
  \item \textbf{Capability} : Permission continue et transférable
  \begin{itemize}
    \item Créé manuellement par votre code
    \item Peut être transféré à d'autres
    \item Utilisé pour des permissions continues (mint, burn, admin, etc.)
  \end{itemize}
\end{itemize}
\end{explainBox}

% ============================================================
\section{Option<T> (Peut-être une Valeur, Peut-être Rien)}

\subsection{Qu'est-ce qu'Option ? (Le Concept)}

Move évite \texttt{null}. Au lieu de cela, il utilise \texttt{Option<T>}.

\begin{explainBox}
\textbf{Le Problème avec null} :
\begin{itemize}
  \item Dans beaucoup de langages, une variable peut être "null" (rien)
  \item Cela cause souvent des bugs (erreur "null pointer")
  \item Move force à être explicite : "Est-ce qu'il y a une valeur ou non ?"
\end{itemize}
\end{explainBox}

\subsubsection{La Structure d'Option}

\[
Option<T> = \text{Some}(T) \ \text{ou} \ \text{None}
\]

\begin{exampleBox}
\textbf{Exemple Concret} :
\begin{itemize}
  \item \texttt{Option<address>} peut être :
  \begin{itemize}
    \item \texttt{Some(0x123...)} → Il y a une adresse
    \item \texttt{None} → Il n'y a pas d'adresse
  \end{itemize}
\end{itemize}
\end{exampleBox}

\subsection{Pourquoi Option est Important}

\begin{explainBox}
\textbf{Cas d'Usage Typiques} :
\begin{itemize}
  \item \textbf{Admin optionnel} : Peut-être qu'il n'y a pas encore d'admin
  \begin{itemize}
    \item Pas d'admin encore → \texttt{None}
    \item Admin existe → \texttt{Some(admin\_address)}
  \end{itemize}
  
  \item \textbf{Valeur par défaut} : Peut-être qu'une valeur n'a pas encore été définie
  \item \textbf{Recherche} : Quand vous cherchez quelque chose qui pourrait ne pas exister
\end{itemize}
\end{explainBox}

\subsection{Exemple Simple}

\begin{lstlisting}[language=Rust, caption=Utilisation Basique d'Option]
use std::option;

// Créer un Option avec une valeur
let a = option::some(10);  // Some(10)

// Créer un Option vide
let b = option::none<u64>();  // None

// Vérifier si c'est Some
if (option::is_some(&a)) {
    // Extraire la valeur (consomme l'Option)
    let x = option::extract(a);  // x = 10
    // Maintenant 'a' n'existe plus (consommé)
}

// Si vous voulez juste lire sans consommer
if (option::is_some(&b)) {
    let y = *option::borrow(&b);  // Emprunte, ne consomme pas
}
\end{lstlisting}

\begin{noteBox}
\textbf{Différence Importante} :
\begin{itemize}
  \item \texttt{extract} : \textbf{Consomme} l'Option (la détruit) et retourne la valeur
  \item \texttt{borrow} : \textbf{Emprunte} l'Option (ne la détruit pas) et retourne une référence
\end{itemize}
\end{noteBox}

\begin{warnBox}
\textbf{⚠️ Erreur Commune} : Si vous utilisez \texttt{extract} sur un \texttt{None}, la transaction échoue. Toujours vérifier avec \texttt{is\_some()} avant d'extraire.
\end{warnBox}

% ============================================================
\section{Vectors (Listes Dynamiques)}

\subsection{Qu'est-ce qu'un Vector ? (Le Concept)}

Un \textbf{vector} est une \textbf{liste dynamique} d'éléments du même type.

\begin{explainBox}
\textbf{Analogie Simple} :
\begin{itemize}
  \item Un vector est comme une \textbf{liste de courses}
  \item Vous pouvez ajouter des éléments
  \item Vous pouvez retirer des éléments
  \item Vous pouvez voir combien d'éléments il y a
  \item Tous les éléments sont du même type (tous des strings, tous des nombres, etc.)
\end{itemize}
\end{explainBox}

\subsubsection{Notation Mathématique}

\[
vector<T> = [T, T, T, \dots]
\]

\subsection{Opérations de Base}

\begin{explainBox}
\textbf{Les Opérations Principales} :
\begin{enumerate}
  \item \textbf{Créer} : \texttt{vector::empty<T>()} - Créer un vector vide
  \item \textbf{Ajouter} : \texttt{vector::push\_back(\&mut v, x)} - Ajouter à la fin
  \item \textbf{Retirer} : \texttt{vector::pop\_back(\&mut v)} - Retirer le dernier
  \item \textbf{Longueur} : \texttt{vector::length(\&v)} - Nombre d'éléments
  \item \textbf{Lire} : \texttt{vector::borrow(\&v, i)} - Lire l'élément à l'index i
\end{enumerate}
\end{explainBox}

\subsection{Exemple Débutant}

\begin{lstlisting}[language=Rust, caption=Exemple Complet d'Utilisation de Vector]
use std::vector;

// 1. Créer un vector vide
let v = vector::empty<u64>();

// 2. Ajouter des éléments
vector::push_back(&mut v, 10);  // v = [10]
vector::push_back(&mut v, 20);  // v = [10, 20]
vector::push_back(&mut v, 30);  // v = [10, 20, 30]

// 3. Voir la longueur
let len = vector::length(&v);  // len = 3

// 4. Lire un élément (sans le retirer)
let first = *vector::borrow(&v, 0);  // first = 10
let second = *vector::borrow(&v, 1); // second = 20

// 5. Retirer le dernier élément
let last = vector::pop_back(&mut v);  // last = 30
// Maintenant v = [10, 20]
\end{lstlisting}

\begin{noteBox}
\textbf{Points Importants} :
\begin{itemize}
  \item \texttt{borrow} retourne une \textbf{référence} - donc vous devez utiliser \texttt{*} pour obtenir la valeur
  \item \texttt{pop\_back} \textbf{retire et retourne} le dernier élément
  \item Les indices commencent à 0 (premier élément = index 0)
\end{itemize}
\end{noteBox}

\begin{warnBox}
\textbf{⚠️ Erreur Commune} : Si vous essayez de \texttt{borrow} avec un index qui n'existe pas (par exemple, index 10 dans un vector de 3 éléments), la transaction échoue. Toujours vérifier la longueur avant d'accéder à un index.
\end{warnBox}

\subsection{Note sur la Performance Sui}

\begin{explainBox}
\textbf{Performance avec Objets Partagés} :
\begin{itemize}
  \item Si un vector est stocké dans un \textbf{objet partagé} et que beaucoup d'utilisateurs le modifient, il y aura des conflits
  \item Cela peut ralentir les transactions
  \item \textbf{Solution commune} : Utiliser plusieurs objets ("sharding") au lieu d'un seul grand vector
\end{itemize}
\end{explainBox}

% ============================================================
\section{Loops (while, loop, break, continue)}

\subsection{Les Loops dans Move (Explicites)}

Move utilise principalement \texttt{while}. Il n'y a pas de boucle \texttt{for} classique.

\begin{explainBox}
\textbf{Pourquoi while ?}
\begin{itemize}
  \item Move force à être explicite sur les conditions
  \item \texttt{while} vous force à gérer manuellement l'index et la condition
  \item Cela rend le code plus clair sur ce qui se passe
\end{itemize}
\end{explainBox}

\subsection{Itérer un Vector}

\begin{lstlisting}[language=Rust, caption=Itérer un Vector avec while]
use std::vector;

let v = vector::empty<u64>();
vector::push_back(&mut v, 10);
vector::push_back(&mut v, 20);
vector::push_back(&mut v, 30);

// Itérer avec while
let i = 0;  // Index de départ
let len = vector::length(&v);  // Longueur du vector

while (i < len) {
    // Lire l'élément à l'index i
    let val = *vector::borrow(&v, i);
    
    // Faire quelque chose avec val
    // (par exemple, l'afficher, le traiter, etc.)
    
    i = i + 1;  // Passer à l'élément suivant
}
\end{lstlisting}

\begin{noteBox}
\textbf{Points Clés} :
\begin{itemize}
  \item Vous devez gérer manuellement l'index \texttt{i}
  \item Vous devez incrémenter \texttt{i} à chaque itération
  \item La condition \texttt{i < len} s'assure que vous ne dépassez pas les limites
\end{itemize}
\end{noteBox}

\subsection{break et continue}

\begin{lstlisting}[language=Rust, caption=Utilisation de break et continue]
let i = 0;
let len = vector::length(&v);

while (i < len) {
    let val = *vector::borrow(&v, i);
    
    // break : Sortir complètement de la boucle
    if (val == 0) {
        break;  // Arrêter la boucle maintenant
    };
    
    // continue : Passer à l'itération suivante
    if (val == 1) {
        i = i + 1;
        continue;  // Ignorer le reste de cette itération
    };
    
    // Code normal ici...
    
    i = i + 1;
}
\end{lstlisting}

\begin{noteBox}
\textbf{Explication} :
\begin{itemize}
  \item \texttt{break} : Sort immédiatement de la boucle (comme "arrêter maintenant")
  \item \texttt{continue} : Passe à l'itération suivante (comme "ignorer le reste et continuer")
  \item N'oubliez pas d'incrémenter \texttt{i} avant \texttt{continue} !
\end{itemize}
\end{noteBox}

\begin{warnBox}
\textbf{⚠️ Attention Gas} : Les loops coûtent du gas. Des loops très grandes peuvent rendre les transactions coûteuses ou échouer à cause des limites de gas. Essayez de garder les loops aussi courtes que possible.
\end{warnBox}

% ============================================================
\section{Events (Comment Votre Contrat Parle au Monde Extérieur)}

\subsection{Qu'est-ce qu'un Event ? (Le Concept)}

Un \textbf{event} est un \textbf{log structuré} émis pendant une transaction.

\begin{explainBox}
\textbf{Analogie avec un Journal} :
\begin{itemize}
  \item Un event est comme une \textbf{entrée dans un journal}
  \item Il enregistre "ce qui s'est passé"
  \item Il ne change pas l'état persistant par lui-même
  \item Il est stocké dans la sortie de transaction (effects)
  \item Les frontends et indexers le lisent
\end{itemize}
\end{explainBox}

\subsubsection{Events vs Objets}

\begin{exampleBox}
\textbf{La Différence} :
\begin{itemize}
  \item \textbf{Objet} : État persistant on-chain (peut être lu plus tard)
  \begin{itemize}
    \item Comme un livre dans une bibliothèque - il reste là
    \item Vous pouvez le lire à tout moment
  \end{itemize}
  
  \item \textbf{Event} : Log de sortie (utile off-chain, pas utilisé comme état)
  \begin{itemize}
    \item Comme une notification - elle dit "quelque chose s'est passé"
    \item Utile pour mettre à jour les UIs, mais pas pour stocker l'état
  \end{itemize}
\end{itemize}
\end{exampleBox}

\subsection{Définir un Event Struct}

Les events ont généralement les abilities \texttt{copy, drop}.

\begin{lstlisting}[language=Rust, caption=Définition d'un Event]
module demo::events {
    use sui::event;
    
    // Un event struct - doit avoir copy et drop
    struct MintEvent has copy, drop {
        owner: address,   // Qui a reçu le mint
        amount: u64       // Combien a été mint
    }
}
\end{lstlisting}

\begin{noteBox}
\textbf{Pourquoi copy et drop ?}
\begin{itemize}
  \item \texttt{copy} : Permet de copier l'event pour le logger
  \item \texttt{drop} : Permet de "jeter" l'event après l'avoir émis
  \item Les events ne sont pas des actifs, donc copy est sûr
\end{itemize}
\end{noteBox}

\subsection{Émettre un Event}

\begin{lstlisting}[language=Rust, caption=Émission d'un Event]
module demo::events {
    use sui::event;
    use sui::tx_context;
    
    struct MintEvent has copy, drop {
        owner: address,
        amount: u64
    }
    
    public fun mint(amount: u64, ctx: &mut TxContext) {
        // ... logique de mint ...
        
        // Émettre l'event
        event::emit(MintEvent {
            owner: tx_context::sender(ctx),  // Qui a appelé la fonction
            amount: amount                    // Combien a été mint
        });
    }
}
\end{lstlisting}

\begin{noteBox}
\textbf{Comment ça fonctionne} :
\begin{itemize}
  \item \texttt{event::emit()} enregistre l'event dans la transaction
  \item Les frontends peuvent écouter ces events
  \item Les indexers peuvent les indexer pour des recherches
  \item C'est excellent pour les mises à jour UI en temps réel
\end{itemize}
\end{noteBox}

\subsection{Events vs Objets (Quand Utiliser Chacun)}

\begin{explainBox}
\textbf{Guide de Décision} :
\begin{itemize}
  \item \textbf{Utilisez un Objet} si :
  \begin{itemize}
    \item Vous avez besoin de lire l'état plus tard
    \item Vous avez besoin de modifier l'état
    \item L'information doit persister sur la blockchain
  \end{itemize}
  
  \item \textbf{Utilisez un Event} si :
  \begin{itemize}
    \item Vous voulez juste notifier que quelque chose s'est passé
    \item Vous voulez que les frontends réagissent
    \item Vous voulez un historique de ce qui s'est passé
    \item L'information n'a pas besoin d'être modifiée plus tard
  \end{itemize}
\end{itemize}
\end{explainBox}

% ============================================================
\section{Tout Mettre Ensemble (Pattern Mini-Projet)}

\subsection{Objectif du Projet}

Nous voulons créer un système avec :
\begin{enumerate}
  \item Un objet \textbf{Config} partagé créé une fois (OTW)
  \item Une \textbf{AdminCap} capability pour le mettre à jour
  \item Des events quand des mises à jour se produisent
  \item Une adresse admin optionnelle (Option)
\end{enumerate}

\subsection{Design Illustratif}

\begin{lstlisting}[language=Rust, caption=Exemple Complet Combinant Tous les Concepts]
module demo::app {
    use sui::object::{UID};
    use sui::tx_context::TxContext;
    use sui::transfer;
    use sui::event;
    use std::option::{Self, Option};

    // OTW pour l'initialisation unique
    struct INIT {}

    // Capability pour les permissions admin
    struct AdminCap has key { 
        id: UID 
    }

    // Objet partagé pour la configuration globale
    struct Config has key {
        id: UID,
        admin: Option<address>,  // Admin optionnel
        threshold: u64
    }

    // Event pour notifier les changements
    struct ConfigUpdated has copy, drop {
        new_threshold: u64
    }

    // Initialiser une fois : créer Config (partagé) et AdminCap (possédé)
    public fun init(w: INIT, ctx: &mut TxContext) {
        // Créer la capability admin
        let cap = AdminCap { id: object::new(ctx) };
        
        // Créer la configuration avec admin optionnel
        let cfg = Config {
            id: object::new(ctx),
            admin: option::some(tx_context::sender(ctx)),  // Some(adresse)
            threshold: 0
        };

        // Donner AdminCap au créateur
        transfer::transfer(cap, tx_context::sender(ctx));

        // Partager Config pour que tout le monde puisse le lire/utiliser
        transfer::share_object(cfg);
    }

    // Seul l'admin peut mettre à jour la config
    public fun set_threshold(
        cap: &AdminCap,  // ✅ Doit avoir AdminCap
        cfg: &mut Config, 
        v: u64
    ) {
        cfg.threshold = v;
        
        // Émettre un event pour notifier
        event::emit(ConfigUpdated { 
            new_threshold: v 
        });
    }
}
\end{lstlisting}

\begin{explainBox}
\textbf{Explication Ligne par Ligne} :
\begin{enumerate}
  \item \textbf{OTW (INIT)} : Garantit que \texttt{init()} ne peut être appelée qu'une fois
  \item \textbf{AdminCap} : Objet capability qui représente la permission admin
  \item \textbf{Config} : Objet partagé contenant la configuration globale
  \item \textbf{Option<address>} : Admin optionnel (peut être Some ou None)
  \item \textbf{ConfigUpdated} : Event émis quand la config change
  \item \textbf{init()} : Crée tout une fois, donne AdminCap au créateur, partage Config
  \item \textbf{set\_threshold()} : Nécessite AdminCap, met à jour Config, émet un event
\end{enumerate}
\end{explainBox}

\begin{warnBox}
\textbf{⚠️ Note} : Ceci est un exemple pédagogique. Dans de vraies applications, vous gérerez soigneusement les règles d'accès et choisirez quels objets sont possédés vs partagés selon vos besoins de performance.
\end{warnBox}

% ============================================================
\section{Erreurs Communes de Débutant (Checklist)}

\begin{enumerate}
  \item \textbf{Donner \texttt{copy} aux actifs} (DANGEREUX !)
  \begin{itemize}
    \item Si un Coin a \texttt{copy}, vous pouvez dupliquer l'argent
    \item \textbf{Solution} : Ne jamais donner \texttt{copy} aux actifs
  \end{itemize}
  
  \item \textbf{Oublier \texttt{id: UID} dans les structs \texttt{has key}}
  \begin{itemize}
    \item Si un struct a \texttt{key}, il DOIT avoir \texttt{id: UID}
    \item \textbf{Solution} : Toujours inclure \texttt{id: UID} dans les objets Sui
  \end{itemize}
  
  \item \textbf{Traiter les events comme du stockage persistant}
  \begin{itemize}
    \item Les events sont des logs, pas de l'état
    \item \textbf{Solution} : Utiliser des objets pour l'état persistant
  \end{itemize}
  
  \item \textbf{Écrire de grandes loops sur des vectors partagés}
  \begin{itemize}
    \item Performance dégradée + gas élevé
    \item \textbf{Solution} : Utiliser plusieurs objets (sharding) ou limiter la taille
  \end{itemize}
  
  \item \textbf{Ne pas utiliser de capabilities pour les actions admin}
  \begin{itemize}
    \item Vérifier seulement \texttt{sender} n'est pas suffisant
    \item \textbf{Solution} : Utiliser des capabilities pour les permissions
  \end{itemize}
  
  \item \textbf{Confondre OTW (init unique) avec Publisher (preuve d'auteur)}
  \begin{itemize}
    \item OTW = initialisation une fois
    \item Publisher = preuve que vous êtes l'auteur
    \item \textbf{Solution} : Comprendre la différence et utiliser chacun correctement
  \end{itemize}
\end{enumerate}

% ============================================================
\section{Modèle Mental Final (Ce qu'il Faut Retenir)}

\begin{center}
\Large\textbf{Sui Move = Objets + Ownership + Abilities + Capabilities}
\end{center}

\subsection{Les Concepts Clés}

\begin{enumerate}
  \item \textbf{Abilities} : Empêchent les comportements non sécurisés (copier l'argent, jeter des actifs)
  
  \item \textbf{Objets} : L'état persistant (avec UID)
  
  \item \textbf{OTW} : Assure l'initialisation unique
  
  \item \textbf{Publisher} : Prouve l'auteur du package et peut protéger les actions admin
  
  \item \textbf{Capabilities} : Les "objets de permission"
  
  \item \textbf{Option} : Donne des valeurs "peut-être" sécurisées (pas de null)
  
  \item \textbf{Vector} : Stocke les listes, les loops itèrent dessus
  
  \item \textbf{Events} : Diffusent ce qui s'est passé au monde off-chain
\end{enumerate}

\subsection{La Philosophie Sui Move}

\begin{noteBox}
\textbf{Le Principe Fondamental} :
\begin{itemize}
  \item \textbf{Sécurité par Design} : Le langage vous force à être explicite sur les actifs
  \item \textbf{Objets comme Permissions} : Posséder un objet = avoir une permission
  \item \textbf{Explicite plutôt qu'Implicite} : Vous devez dire exactement ce que vous voulez faire
  \item \textbf{Composition} : Construisez des systèmes complexes à partir de briques simples
\end{itemize}
\end{noteBox}

\end{document}
